\documentclass[11pt,reqno]{amsart}
\usepackage[margin=1.15in]{geometry}
\usepackage{amsmath,amssymb,amsthm,mathtools}
\usepackage{graphicx}
\usepackage{booktabs}
\usepackage[colorlinks=true,linkcolor=blue,citecolor=blue,urlcolor=blue]{hyperref}
\usepackage{microtype}

\theoremstyle{plain}
\newtheorem{theorem}{Theorem}
\newtheorem{proposition}[theorem]{Proposition}
\newtheorem{corollary}[theorem]{Corollary}
\newtheorem{lemma}[theorem]{Lemma}
\theoremstyle{definition}
\newtheorem{definition}[theorem]{Definition}
\newtheorem{remark}[theorem]{Remark}
\newtheorem{observation}[theorem]{Observation}

\DeclareMathOperator{\sqf}{sqf}
\DeclareMathOperator{\ord}{ord}
\newcommand{\Art}{\mathrm{A}}
\newcommand{\Z}{\mathbb{Z}}
\newcommand{\Q}{\mathbb{Q}}
\newcommand{\legs}[2]{\left(\tfrac{#1}{#2}\right)}

\begin{document}

\title[The layers of the consecutive-prime Artin correlation]
{The layers of the consecutive-prime Artin correlation:\\
only the abelian one correlates}

\author{Josh Bald}
\address{Ontario, Canada}
\email{jpbald93@gmail.com}

\subjclass[2020]{Primary 11A07; Secondary 11N05, 11R32, 11Y16}
\keywords{Artin's conjecture, primitive roots, consecutive primes, power
residues, Kummer extensions, Lemke Oliver--Soundararajan bias}

\date{\today}

\begin{abstract}
The primitive-root (``Artin'') condition for a base $a$ at a prime $p$
factors into layers: for each prime $\ell \mid p - 1$, the element $a$ must
not be an $\ell$-th power residue. Earlier papers in this series measured a
systematic correlation between the Artin statuses of consecutive primes and
traced it, empirically, to the quadratic layer and to the joint distribution
of $p - 1$'s factorizations. This paper closes the remaining question: do the
higher layers ($\ell \ge 3$) contribute correlation of their own? We show
that they cannot contribute through gap classes --- for $\ell \ge 3$ the
$\ell$-th power residue symbol of a fixed base is not determined by any
congruence condition on $p$, because $\Q(\zeta_\ell, a^{1/\ell})$ is
non-abelian over $\Q$, so no analogue of the quadratic exclusion laws exists
--- and we verify computationally that they do not contribute at all: over
all consecutive prime pairs below $10^9$ where the $\ell$-th condition binds
at both primes, the $\ell$-th symbol statuses are statistically independent
for every $\ell \in \{3,5,7,11,13\}$ and every base tested, while the
\emph{binding events} $\ell \mid p_n - 1$ and $\ell \mid p_{n+1} - 1$ are
strongly anticorrelated (e.g.\ $\phi = -0.133$, $z = -150$, for $\ell = 3$ at
$2\times 10^7$) --- the Lemke Oliver--Soundararajan bias modulo $\ell$. The
consecutive-prime Artin correlation is therefore carried entirely by the
abelian ($\ell = 2$) layer and by the binding structure of $p-1$: a
dichotomy aligned exactly with the abelian/non-abelian dichotomy of the
Kummer extensions. As a corollary, the conditional model of the sixth paper
in this series, which uses only quadratic and binding data, is not an
approximation that happens to work: no further layer terms exist to add.
\end{abstract}

\maketitle

\section{Introduction}
\label{sec:intro}

Let $a \ge 2$ be a non-square integer. A prime $p \nmid a$ has $a$ as a
primitive root if and only if, for every prime $\ell \mid p - 1$, the element
$a$ is not an $\ell$-th power residue modulo $p$:
\begin{equation}
\label{eq:layers}
  p \in \Art_a
  \iff
  a^{(p-1)/\ell} \not\equiv 1 \pmod p
  \quad \text{for every prime } \ell \mid p - 1 .
\end{equation}
We refer to the condition at $\ell$ as the \emph{$\ell$-th layer}. The
layer at $\ell$ \emph{binds} at $p$ when $\ell \mid p - 1$; the layer at
$\ell = 2$ binds at every odd prime, and its condition is exactly
$\chi_a(p) = -1$ for the quadratic character $\chi_a = \legs{\sqf(a)}{\cdot}$.

Earlier papers in this series measured the correlation
\[
  \delta(a; x) = P\bigl(p_{n+1} \in \Art_a \mid p_n \in \Art_a\bigr)
               - P\bigl(p_{n+1} \in \Art_a \mid p_n \notin \Art_a\bigr)
\]
over consecutive primes up to $x = 10^9$ and found it nonzero for every base
tested ($|z|$ up to $479$), organised by the conductor of $\chi_a$, with
unconditional \emph{exclusion laws} in the quadratic layer: gap classes
modulo the conductor in which two consecutive primes can never both be Artin.
A companion paper constructed a conditional Hardy--Littlewood model from the
quadratic layer and the factorization structure of $p-1$ alone, and found it
reproduces $\delta(a)$ to within $4\%$ for twelve bases, with channel-level
$R^2 > 0.999$.

That success poses a question the model cannot answer for itself: \emph{why
is nothing missing?} The Artin condition~\eqref{eq:layers} involves every
layer. If the cubic layer, say, carried its own consecutive-prime
correlation --- an analogue of the quadratic exclusion laws, or any gap-class
structure --- the model would be incomplete and its agreement partly
accidental. The present paper shows that the higher layers can carry no such
structure, in two steps:

\begin{itemize}
\item \textbf{No exclusion laws above $\ell = 2$}
      (Section~\ref{sec:noexclusion}). The quadratic exclusion laws exist
      because $\chi_a(p)$ is periodic in $p$. For $\ell \ge 3$ the $\ell$-th
      power residue status of a fixed integer $a$ is \emph{not} a congruence
      condition on $p$: the splitting field $\Q(\zeta_\ell, a^{1/\ell})$ is
      non-abelian over $\Q$, and a Chebotarev-style argument shows its
      splitting behaviour cannot be captured by residue classes. Gauss's
      classical criterion --- $2$ is a cubic residue mod $p$ iff
      $p = L^2 + 27M^2$ --- is the concrete face of this obstruction.
\item \textbf{Empirically, no correlation at all}
      (Section~\ref{sec:measurement}). Across all consecutive prime pairs
      below $10^9$ at which the $\ell$-th layer binds at both primes, the
      symbol statuses are statistically independent, for every
      $\ell \in \{3, 5, 7, 11, 13\}$ and every base
      $a \in \{2,3,5,6,7,10,11,13,15\}$. The \emph{binding events}
      themselves, by contrast, are strongly anticorrelated --- pure Lemke
      Oliver--Soundararajan bias modulo $\ell$ \cite{LOS2016}.
\end{itemize}

The conclusion (Section~\ref{sec:discussion}) is a dichotomy that matches
the Galois theory exactly: \emph{the abelian layer correlates; the
non-abelian layers are independent}. All consecutive-prime structure in the
Artin property flows through (i) the quadratic character and (ii) which
layers bind --- both of which are congruence data --- and none through the
non-abelian symbols. In particular the companion model's completeness is
structural, not fortuitous.

\section{No exclusion laws above the quadratic layer}
\label{sec:noexclusion}

The quadratic exclusion laws of the second paper rest on one fact: $\chi_a$
is a Dirichlet character modulo the conductor $f$, so $\chi_a(p)$ depends
only on $p \bmod f$, and a gap class that forces $\chi_a(q) = -\chi_a(p)$
forces one of the two primes out of $\Art_a$. The following proposition
shows this mechanism is exclusive to $\ell = 2$.

\begin{proposition}
\label{prop:nonabelian}
Let $\ell \ge 3$ be prime and let $a \ge 2$ be an integer that is not a
perfect $\ell$-th power. There is no modulus $M$ and set
$S \subseteq (\Z/M\Z)^\times$ such that, for all sufficiently large primes
$p \equiv 1 \pmod \ell$,
\[
  a \text{ is an $\ell$-th power residue mod } p
  \iff p \bmod M \in S .
\]
\end{proposition}

\begin{proof}
Let $K = \Q(\zeta_\ell, a^{1/\ell})$. For a prime
$p \equiv 1 \pmod \ell$ not dividing $a\ell$, the condition ``$a$ is an
$\ell$-th power residue mod $p$'' is equivalent to $p$ splitting completely
in $K$ (standard Kummer theory: $p$ splits completely in $\Q(\zeta_\ell)$
because $p \equiv 1 \pmod \ell$, and then splits completely in $K$ iff
$x^\ell - a$ has a root --- equivalently all roots --- mod $p$).

Suppose such $M$ and $S$ existed. Then the set of rational primes splitting
completely in $K$ would be, up to finitely many exceptions, a union of
arithmetic progressions modulo $\mathrm{lcm}(M, \ell)$. A number field whose
set of completely-split primes is (up to density-zero exceptions) a union of
arithmetic progressions mod $N$ has the same completely-split set as some
subfield of $\Q(\zeta_N)$; since the completely-split set determines the
Galois closure (a standard consequence of the Chebotarev density theorem),
the Galois closure of $K$ would be contained in $\Q(\zeta_N)$ and hence
\emph{abelian} over $\Q$.

But $K$ is its own Galois closure ($\zeta_\ell \in K$), with
$\mathrm{Gal}(K/\Q) \cong \Z/\ell\Z \rtimes (\Z/\ell\Z)^\times$ of order
$\ell(\ell-1)$: it is non-abelian for every $\ell \ge 3$, because the
subgroup fixing $\Q(\zeta_\ell)$ (translations $a^{1/\ell} \mapsto
\zeta_\ell^j a^{1/\ell}$) is normal and acted on nontrivially by
$(\Z/\ell\Z)^\times$, which has order $\ell - 1 \ge 2$. (Non-degeneracy of
the action requires $a$ not an $\ell$-th power, which we assumed; the
semidirect product is then genuinely non-abelian.) Contradiction.
\end{proof}

\begin{corollary}
\label{cor:noexclusion}
For $\ell \ge 3$ there is no gap-class exclusion law at the $\ell$-th layer:
no modulus $M$ and residue class $g_0$ such that consecutive primes
$p, q = p + g$ with $g \equiv g_0 \pmod M$, both $\equiv 1 \pmod \ell$, are
forbidden from both having $a$ as an $\ell$-th power non-residue (or any
other fixed pair of symbol statuses).
\end{corollary}

\begin{proof}
A gap-class law forcing symbol statuses would, applied to the pairs
$(p, p+g)$ ranging over all admissible $p$, determine the symbol status of
$p + g$ from congruence data of $p$ and $g$ alone --- in particular the
symbol status at a prime $q$ would be a congruence condition on $q$
(intersect over the available $p$; the LOS equidistribution of $p$ over
admissible classes makes every class available). This contradicts
Proposition~\ref{prop:nonabelian}.
\end{proof}

\begin{remark}
The classical criteria make the obstruction tangible. $2$ is a cubic residue
modulo $p \equiv 1 \pmod 3$ iff $p = L^2 + 27M^2$ (Gauss); $2$ is a quartic
residue modulo $p \equiv 1 \pmod 8$ iff $p = L^2 + 64M^2$. Representation by
a quadratic form of discriminant $-108$ or $-256$ is class-field data over
$\Q(\zeta_3)$ resp.\ $\Q(i)$ --- abelian over the quadratic field, but not
over $\Q$ --- and is well known not to reduce to congruences in $p$.
\end{remark}

\begin{remark}
Proposition~\ref{prop:nonabelian} does not by itself force statistical
independence of the symbols at consecutive primes; it only removes the
congruence mechanism that produces exclusion laws and, more generally, any
correlation \emph{mediated by joint residue classes} of
$(p_n, p_{n+1}, g)$ --- which is precisely the mechanism (Hardy--Littlewood
singular series, LOS bias) known to govern consecutive-prime dependence.
Whether some other mechanism correlates the non-abelian symbols across
consecutive primes is an empirical question, which we now settle at scale.
\end{remark}

\section{Measurement: the higher symbols are independent}
\label{sec:measurement}

\subsection{Setup}
For each prime $\ell \in \{3,5,7,11,13\}$ and base
$a \in \{2,3,5,6,7,10,11,13,15\}$, define at each prime $p$ with
$\ell \mid p - 1$ the symbol status
$s_\ell(a, p) = 1$ if $a^{(p-1)/\ell} \equiv 1 \pmod p$, else $0$. We
accumulate, over consecutive prime pairs $(p_n, p_{n+1})$ below $10^9$:
(i) the $2\times2$ table of binding events
$(\ell \mid p_n - 1,\ \ell \mid p_{n+1} - 1)$;
(ii) for pairs binding at both primes, the $2\times2$ table of
$(s_\ell(a, p_n),\ s_\ell(a, p_{n+1}))$;
(iii) for the $\ell = 2$ layer, the table of quadratic-character signs; and
(iv) the Artin-status table, as a consistency anchor to the earlier papers.
We report the $\phi$ coefficient (Pearson correlation of the two indicators)
and $z = \phi\sqrt{n}$.

The computation factors each $p - 1$ once and evaluates all layers and bases
in a single pass (the scheme of the second paper); a full run to $10^9$
takes about an hour on one core. The pilot stage was implemented
independently in \texttt{sympy} and agrees with the C implementation to all
reported digits at $2\times10^7$.

\subsection{Results}

\begin{table}[t]
\centering
\caption{Binding correlations at $x = 10^9$: $\phi$ for the events
$\ell \mid p_n - 1$ versus $\ell \mid p_{n+1} - 1$ over all consecutive
pairs. All entries are exact counts over $50{,}847{,}527$ pairs.}
\label{tab:bind}
\begin{tabular}{rrr}
\toprule
$\ell$ & $\phi$ & $z$ \\
\midrule
$3$  & $-0.10977$  & $-782.7$ \\
$5$  & $-0.08909$  & $-635.2$ \\
$7$  & $-0.07884$  & $-562.2$ \\
$11$ & $-0.05954$ & $-424.6$ \\
$13$ & $-0.05763$ & $-411.0$ \\
\bottomrule
\end{tabular}
\end{table}

\begin{table}[t]
\centering
\caption{Symbol correlations at $x = 10^9$, conditional on binding at both
primes: $\phi$ (with $z$ in parentheses) for
$s_\ell(a, p_n)$ versus $s_\ell(a, p_{n+1})$.
Maximum $|z|$ over all 45 combinations: $2.4$, consistent with 45 null tests.}
\label{tab:sym}
\begin{tabular}{rrrrrr}
\toprule
$a$ & $\ell=3$ & $\ell=5$ & $\ell=7$ & $\ell=11$ & $\ell=13$ \\
\midrule
$2$ & $-0.00024$ ($-0.8$) & $+0.00080$ ($+1.2$) & $-0.00034$ ($-0.3$) & $-0.00028$ ($-0.1$) & $-0.00077$ ($-0.3$) \\
$3$ & $+0.00015$ ($+0.5$) & $+0.00047$ ($+0.7$) & $+0.00150$ ($+1.4$) & $-0.00137$ ($-0.7$) & $-0.00020$ ($-0.1$) \\
$5$ & $-0.00008$ ($-0.3$) & $-0.00140$ ($-2.1$) & $+0.00035$ ($+0.3$) & $+0.00155$ ($+0.8$) & $+0.00192$ ($+0.7$) \\
$6$ & $-0.00002$ ($-0.1$) & $+0.00031$ ($+0.5$) & $+0.00000$ ($+0.0$) & $+0.00413$ ($+2.0$) & $+0.00141$ ($+0.5$) \\
$7$ & $-0.00026$ ($-0.9$) & $+0.00063$ ($+1.0$) & $+0.00045$ ($+0.4$) & $+0.00179$ ($+0.9$) & $+0.00061$ ($+0.2$) \\
$10$ & $+0.00023$ ($+0.8$) & $+0.00080$ ($+1.2$) & $+0.00237$ ($+2.2$) & $+0.00138$ ($+0.7$) & $-0.00051$ ($-0.2$) \\
$11$ & $-0.00031$ ($-1.0$) & $+0.00017$ ($+0.3$) & $-0.00174$ ($-1.6$) & $-0.00418$ ($-2.0$) & $-0.00302$ ($-1.1$) \\
$13$ & $-0.00012$ ($-0.4$) & $+0.00011$ ($+0.2$) & $-0.00126$ ($-1.2$) & $-0.00084$ ($-0.4$) & $-0.00023$ ($-0.1$) \\
$15$ & $-0.00027$ ($-0.9$) & $-0.00002$ ($-0.0$) & $-0.00124$ ($-1.1$) & $-0.00207$ ($-1.0$) & $+0.00668$ ($+2.4$) \\
\bottomrule
\end{tabular}
\end{table}

Two facts, in deliberate contrast:

\begin{observation}[Binding layer: strong anticorrelation]
\label{obs:bind}
The binding events are anticorrelated for every $\ell$, at overwhelming
significance. This is the Lemke Oliver--Soundararajan bias modulo $\ell$ in
its purest form: a prime $\equiv 1 \pmod \ell$ repels a successor in its own
residue class, and $1 \bmod \ell$ is a single class among $\ell - 1$.
\end{observation}

\begin{observation}[Symbol layer: independence]
\label{obs:sym}
Conditional on binding at both primes, the symbol statuses are independent:
across all $45$ (base, $\ell$) combinations, no $|z|$ exceeds the range
expected from $45$ null tests. The marginal symbol frequencies are
$1/\ell$ to high accuracy, as they must be (Kummer--Chebotarev
equidistribution), and the joint frequencies factor.
\end{observation}

\subsection{Consistency anchor}
The same run reproduces the Artin-status correlations $\delta(a)$ of the
earlier papers to the last digit for all nine bases
(e.g.\ $\delta(2) = -0.050199$, $\delta(15) = +0.012196$, matching to six decimals), confirming that the layer data and the Artin data come from
the identical prime stream.

\section{Discussion}
\label{sec:discussion}

\subsection*{The dichotomy}
Assemble the three findings: (i) the quadratic symbol is periodic in $p$,
inherits the consecutive-prime residue bias, and produces exclusion laws and
the conductor-organised correlation of the earlier papers; (ii) the binding
events at every $\ell$ are congruence data ($p \bmod \ell$) and are strongly
anticorrelated; (iii) the non-abelian symbols at every $\ell \ge 3$ carry no
consecutive-prime correlation at all. The split is exactly the
abelian/non-abelian split of the splitting fields: $\Q(\sqrt{\sqf a})$ and
$\Q(\zeta_\ell)$ are abelian over $\Q$, so their splitting data are
congruence conditions, and consecutive primes are biased in congruence data;
$\Q(\zeta_\ell, a^{1/\ell})$ for $\ell \ge 3$ is non-abelian, its splitting
data are invisible to residue classes, and consecutive primes appear to
carry no bias in them whatsoever.

We find the cleanliness of (iii) worth emphasising. Nothing in
Proposition~\ref{prop:nonabelian} \emph{forbids} a correlation of
non-abelian symbols --- it only rules out the congruence mechanism. The
Hardy--Littlewood/LOS framework, as far as it is understood, biases exactly
the joint congruence data of consecutive primes; our measurement is
consistent with the sharpest reading of that framework: \emph{consecutive
primes are dependent only through congruence data}. Symbol independence at
$45$ combinations, with the binding channels simultaneously showing
$|z| > 100$ in the same data stream, is a strong internal control: the
pipeline demonstrably detects correlation where correlation exists.

\subsection*{Consequences for the series}
First, the conditional model of the sixth paper is structurally complete:
its inputs (quadratic layer, binding structure, gap classes) exhaust the
channels through which consecutive-prime dependence can enter the Artin
property, so its $R^2 > 0.999$ channel agreement is not a fortunate
truncation. Second, the ``cubic or quartic sharing'' speculation recorded in
the third paper as a possible source of a small residual floor is now
implausible: the cubic symbols themselves are uncorrelated, so any residual
in that paper's cross-base analysis must come from the binding structure or
from finite-size effects, not from a cubic analogue of quadratic
entanglement. Third, the exclusion-law programme of the second paper is
complete as stated: there are no higher-layer exclusion laws to find.

\subsection*{Limitations}
The independence claim is statistical, bounded by the resolution of the
census ($\approx 6\times10^{-4}$ in $\phi$ at $z = 4$ for $\ell = 3$;
weaker for larger $\ell$, where binding pairs are scarcer). A correlation of
order $10^{-4}$ at $\ell = 3$ would evade us. We see no mechanism that would
produce one, but the bound is what the data give. Second,
Proposition~\ref{prop:nonabelian} is stated for prime $\ell$ and non-$\ell$-th-power
$a$; composite indices add nothing new (the layers of \eqref{eq:layers} are
prime), but bases that are perfect $\ell$-th powers have degenerate layers
and are excluded throughout, as in the earlier papers.

\section*{Acknowledgements}
The author thanks the developers of the open-source tools used in this work
(GCC, Python, \texttt{sympy}, \texttt{matplotlib}).

\subsection*{Use of generative AI}
The author used Anthropic's Claude (models Claude Opus 4.5 and Claude
Opus 5) as an assistant: for drafting and revising exposition, for writing
and debugging the sieve and analysis programs, and for exploratory
discussion. It is recorded, in the interest of transparency, that this paper
began as a search for \emph{cubic exclusion laws}; the structural obstruction
of Proposition~\ref{prop:nonabelian} and the null measurement replaced that
plan --- the third time in this series that an initially proposed mechanism
was corrected by the data. All mathematical claims were verified
computationally from exact counts; the proofs were checked by hand; the
author is solely responsible for the content.

\section*{Declaration of interest}
The author reports there are no competing interests to declare.

\section*{Funding}
This research received no external funding.

\section*{Data availability}
All code and results are publicly available at
\url{https://github.com/jpbald93/consecutive-artin}, including the one-pass
layer census (C), the independent \texttt{sympy} pilot, the exact
contingency tables for every (base, $\ell$) combination, and the scripts
reproducing every table.

\begin{thebibliography}{99}

\bibitem{Gauss}
C.~F. Gauss,
\emph{Disquisitiones Arithmeticae},
Leipzig, 1801; \S358 (cubic), see also the criterion for biquadratic residues
in the Comm. Soc. Reg. Sci. G\"ottingen (1828/1832).

\bibitem{Hooley1967}
C.~Hooley,
\emph{On Artin's conjecture},
J. Reine Angew. Math. \textbf{225} (1967), 209--220.

\bibitem{Lemmermeyer}
F.~Lemmermeyer,
\emph{Reciprocity Laws: From Euler to Eisenstein},
Springer Monographs in Mathematics, Springer-Verlag, Berlin, 2000.

\bibitem{LOS2016}
R.~J. Lemke~Oliver and K.~Soundararajan,
\emph{Unexpected biases in the distribution of consecutive primes},
Proc. Natl. Acad. Sci. USA \textbf{113} (2016), no.~31, E4446--E4454.

\bibitem{Lenstra1977}
H.~W. Lenstra, Jr.,
\emph{On Artin's conjecture and Euclid's algorithm in global fields},
Invent. Math. \textbf{42} (1977), 201--224.

\bibitem{Moree2012}
P.~Moree,
\emph{Artin's primitive root conjecture --- a survey},
Integers \textbf{12A} (2012), Article A13, 100~pp.

\end{thebibliography}

\end{document}
